%%%%%%%%%%%%%%%%%%%Pour les index%%%%%%%%%%%%%%%%
%%%%%%%%Gestion des index pour le document%%%%%%%
%%%%%%%%%%%%avec éléments d'évaluation%%%%%%%%%%%
\makeindex[name=q, title=Index des questions selon leur type\label{indexq}, options={-M styles/questions.xdy -L french -C utf8}]
\makeindex[name=a, title=Index des auteurs\label{indexa},%
options={-M styles/perso.xdy -L french -C utf8}]
\makeindex[name=p, title=Index des questions selon les axes du programme\label{indexp},%
options={-M styles/questions.xdy -L french -C utf8}]
\makeindex[name=g, title=Index des textes selon leur genre\label{indexg},%
options={-M styles/genre.xdy -L french -C utf8}]
\makeindex[name=n, title=Index des sujets par lieux et années\label{indexn},%
options={-M styles/numeros.xdy -L french -C utf8}]
\makeindex[name=t, title=Index thématique\label{indext}, options={-M styles/themes.xdy}]
\makeindex[name=w, title=Index des questions selon les axes du programme puis selon leur type\label{indexw}, options={-M styles/themesquestion.xdy}]
\makeindex[name=v, title=Index des questions selon les types de question puis selon les axes du programme\label{indexv}, options={-M styles/questionthemes.xdy -L french -C utf8}]%questionthemes.ist
%\makeindex[name=c, title={Index des questions avec des éléments d'évaluation\label{indexc}}, options={-M styles/corrige.xdy}]%corrige.ist
%\makeindex[name=r, title={Index des références culturelles, œuvres et auteurs cités dans les éléments d'évaluation\label{indexac}}, options={-M styles/corrigeauteurs.xdy -L french -C utf8}]

%%%Début de la configuration de xkeyval%%%
\makeatletter
%%%%%%%%%%%%%%%%%%%Pour les sujets%%%%%%%%%%%%%%%%

%% 1. Définition des clés pour les sujets 
\define@cmdkey	[SUJ]	{sujet}	{semestre}		{}
\define@cmdkey	[SUJ]	{sujet}	{lieu}			{}
\define@cmdkey	[SUJ]	{sujet}	{annee}			{}
\define@cmdkey	[SUJ]	{sujet}	{type}			{}
\define@cmdkey	[SUJ]	{sujet}	{jour}			{}
\define@cmdkey	[SUJ]	{sujet}	{nosujet}		{}
\define@cmdkey	[SUJ]	{sujet}	{prenomauteur}	{}
\define@cmdkey	[SUJ]	{sujet}	{initialeprenom}{}
\define@cmdkey	[SUJ]	{sujet}	{particule}		{}
\define@cmdkey	[SUJ]	{sujet}	{nomauteur}		{}
\define@cmdkey	[SUJ]	{sujet}	{titre}			{}
\define@cmdkey	[SUJ]	{sujet}	{datepub}		{}
\define@cmdkey	[SUJ]	{sujet}	{ajoutref}		{}
\define@cmdkey	[SUJ]	{sujet}	{genre}			{}
\define@cmdkey	[SUJ]	{sujet}	{corrige}		{}
\define@cmdkey	[SUJ]	{sujet}	{liencorrige}	{}
%\define@cmdkey	[SUJ]	{sujet}	{texte}			{}
\define@cmdkey	[SUJ]	{sujet}	{typequn}		{}
\define@cmdkey	[SUJ]	{sujet}	{typeqdeux}		{}
\define@cmdkey	[SUJ]	{sujet}	{qun}			{}
\define@cmdkey	[SUJ]	{sujet}	{quncsv}		{}
\define@cmdkey	[SUJ]	{sujet}	{qdeux}			{}
\define@cmdkey	[SUJ]	{sujet}	{qdeuxcsv}		{}
\define@cmdkey	[SUJ]	{sujet}	{themeprogrammequn}	{}
\define@cmdkey	[SUJ]	{sujet}	{themeprogrammeqdeux}	{}

%% 2. Définition des commandes pour accéder aux valeurs des clés
\newcommand{\semestre}{\cmdSUJ@sujet@semestre}
\newcommand{\lieu}{\cmdSUJ@sujet@lieu}
\newcommand{\annee}{\cmdSUJ@sujet@annee}
\newcommand{\type}{\cmdSUJ@sujet@type}
\newcommand{\jour}{\cmdSUJ@sujet@jour}
\newcommand{\prenomauteur}{\cmdSUJ@sujet@prenomauteur}
\newcommand{\initialeprenom}{\cmdSUJ@sujet@initialeprenom}
\newcommand{\particule}{\cmdSUJ@sujet@particule}
\newcommand{\nomauteur}{\cmdSUJ@sujet@nomauteur}
\newcommand{\titre}{\cmdSUJ@sujet@titre}
\newcommand{\datepub}{\cmdSUJ@sujet@datepub}
\newcommand{\ajoutref}{\cmdSUJ@sujet@ajoutref}
\newcommand{\genre}{\cmdSUJ@sujet@genre}
\newcommand{\corrige}{\cmdSUJ@sujet@corrige}
\newcommand{\liencorrige}{\cmdSUJ@sujet@liencorrige}
%\newcommand{\texte}{\cmdSUJ@sujet@texte}
\newcommand{\typequn}{\cmdSUJ@sujet@typequn}
\newcommand{\typeqdeux}{\cmdSUJ@sujet@typeqdeux}
\newcommand{\qun}{\cmdSUJ@sujet@qun}
\newcommand{\qdeux}{\cmdSUJ@sujet@qdeux}
\newcommand{\themeprogrammequn}{\cmdSUJ@sujet@themeprogrammequn}
\newcommand{\themeprogrammeqdeux}{\cmdSUJ@sujet@themeprogrammeqdeux}

%% 3. Valeurs par défaut des clés 
\presetkeys[SUJ]{sujet}{%
				nomauteur = Anonyme,
				prenomauteur =,
				initialeprenom =,
				particule =,
				titre = {LeTitre},
				semestre = HQRS,
				lieu = Métropole,
				annee = 2020,
				type =,
				jour = 0,
				datepub = 4545,
				ajoutref =,
				genre = \genrenondefini{},
				corrige =,
				liencorrige=,
				typequn = \litt,
				typeqdeux = \phil,
				qun =,
				quncsv=,
				qdeux =,
				qdeuxcsv=,
				themeprogrammequn=,
				themeprogrammeqdeux=,}{}
%				texte = {Lorem ipsum dolor sit amet},}
%%%Fin de la configuration de xkeyval

% 4. Macros utilitaires pour le formatage et l'affichage
% 4.1 Formatage du nom de l'auteur selon les cas
\newcommand{\formatauteur}{%
	\ifthenelse{\equal{\particule}{}}%
	{\initialeprenom.~\textsc{\nomauteur}}%
	{\IfEndWith{\particule}{'}%
	{\initialeprenom.~\textsc{\particule\nomauteur}}%
	{\initialeprenom.~\textsc{\particule{}~\nomauteur}}%
	}%
	}

% 4.2 Formatage de l'en-tête
\newcommand{\formatsessionentete}{%
    \ifthenelse{\equal{\type}{}}%
    {\lieu{}~\annee{} (Jour~\jour{})~:~}% Sans type spécifié
    {\lieu{}~\annee{} (\type, Jour~\jour{})~:~}% Avec type spécifié
}

%% 5. Macros pour l'indexation et les tables des matières
% 5.0 Commandes variables (visant à être écrasées à chaque fois que la fonction sera exécutée sur chaque texte)
% 5.0.1. Variables pour le format de l'auteur
\providecommand{\tocauteurformat}{}
\providecommand{\indexauteurformat}{}

% 5.0.2. Variables pour le format des informations temporelles
\providecommand{\tocanneformat}{}
\providecommand{\indexanneformat}{}

% 5.0.3. Variable pour la clé de tri de l'index
\providecommand{\indexcletri}{}

% 5.1 Gestion de l'index des questions et des compteurs
% 5.1.1 Indexation de l'auteur
\newcommand{\indexerauteur}{%
    \ifthenelse{\equal{\prenomauteur}{}}%
    {\index[a]{\nomauteur{}@\textsc{\nomauteur}}}% Indexer sans prénom
	    {\ifthenelse{\equal{\particule}{}}%Test de particule
	    	{\index[a]{\nomauteur,~\prenomauteur{}@\textsc{\nomauteur}~\prenomauteur{}}}%Sans particule
	    	{\IfEndWith{\particule}{'}%Test de la fin de particule
	    	{\index[a]{\nomauteur,~\prenomauteur{}@\textsc{\particule\nomauteur{}},~\prenomauteur}}
	    	{\index[a]{\nomauteur,~\prenomauteur{}@\textsc{\particule{}~\nomauteur{}},~\prenomauteur}}
	    	}
	    }
	}%Fin de la commande


% 5.1.2 Index et mise en page des questions
\newcommand{\indexerquestions}{%
	\begin{envquestion}
	
	% ==========================================
	% 1. AFFICHAGE ET INDEXATION DE BASE [q], [p] et [w]
	% ==========================================
	\noindent\question{Question d'interprétation \typequn~:}\par\nopagebreak
	\noindent\qun\index[q]{Question d'interprétation \typequn!\qun}
	
	\spcsujet
	
	\noindent\question{Essai \typeqdeux~:}\par\nopagebreak
	\noindent\qdeux\index[q]{Essai \typeqdeux!\qdeux}
	
	% Index [p] (Thème ! Question)
	\index[p]{{\csname\themeprogrammequn\endcsname}!\qun}
	\index[p]{{\csname\themeprogrammeqdeux\endcsname}!\qdeux}
	
	% Index [w] (Thème ! Type ! Question) -> Plus besoin de conditions ici !
	\index[w]{\csname\themeprogrammequn\endcsname!Question d'interprétation \typequn!\qun}
	\index[w]{\csname\themeprogrammeqdeux\endcsname!Essai \typeqdeux!\qdeux}

	% ==========================================
	% 2. INDEX [v] POUR LA QUESTION 1 (INTERPRÉTATION)
	% ==========================================
	\ifthenelse{\equal{\typequn}{\litt}}%
		{% --- Q1 est Littéraire (Ancre : intlitt...) ---
		 \index[v]{Question d'interprétation \typequn @ Question d'interprétation \typequn !%
		           \csname\themeprogrammequn\endcsname @ \protect\hypertarget{intlitteraire\themeprogrammequn}{\csname\themeprogrammequn\endcsname} !%
		           \qun}%
		}
		{% --- Q1 est Philosophique (Ancre : intphil...) ---
		 \index[v]{Question d'interprétation \typequn @ Question d'interprétation \typequn !%
		           \csname\themeprogrammequn\endcsname @ \protect\hypertarget{intphilosophique\themeprogrammequn}{\csname\themeprogrammequn\endcsname} !%
		           \qun}%
		}

	% ==========================================
	% 3. INDEX [v] POUR LA QUESTION 2 (ESSAI)
	% ==========================================
	\ifthenelse{\equal{\typeqdeux}{\litt}}%
		{% --- Q2 est Littéraire (Ancre : esslitt...) ---
		 \index[v]{Essai \typeqdeux @ Essai \typeqdeux !%
		           \csname\themeprogrammeqdeux\endcsname @ \protect\hypertarget{reflitteraire\themeprogrammeqdeux}{\csname\themeprogrammeqdeux\endcsname} !%
		           \qdeux}%
		}
		{% --- Q2 est Philosophique (Ancre : essphil...) ---
		 \index[v]{Essai \typeqdeux @ Essai \typeqdeux !%
		           \csname\themeprogrammeqdeux\endcsname @ \protect\hypertarget{refphilosophique\themeprogrammeqdeux}{\csname\themeprogrammeqdeux\endcsname} !%
		           \qdeux}%
		}
	\end{envquestion}
}
	
%5.1.3 Gestion des compteurs
\newcommand{\inter}{int}
\newcommand{\refle}{ref}
\newcommand{\litte}{litteraire}
%\newcommand{\philo}{philosophique}
\newcommand{\intlitt}{intlitteraire}
\newcommand{\intphil}{intphilosophique}
\newcommand{\reflitt}{reflitteraire}
\newcommand{\refphil}{refphilosophique}

\newcommand{\compteurquestions}{
	\stepcounter{sujetint\themeprogrammequn}%
	\stepcounter{sujetref\themeprogrammeqdeux}%
	\stepcounter{sujet\themeprogrammequn}%
	\stepcounter{sujet\themeprogrammeqdeux}%
	\ifthenelse{\equal{\typequn}{\litt}}%
	%%%Si l'interprétation est littéraire, on augmente l'interprétation litt. liés et la réf. philosophique ; ainsi que les compteurs liés.
	{\stepcounter{sujetintlitteraire\themeprogrammequn}
	\stepcounter{nosujetintlitteraire}
	\stepcounter{sujetrefphilosophique\themeprogrammeqdeux}
	\stepcounter{nosujetrefphilosophique}%
	}
	{%%%Sinon, on augmente l'interprétation phil. et la réf. litt. ; ainsi que les compteurs liés.
	\stepcounter{sujetintphilosophique\themeprogrammequn}
	\stepcounter{nosujetintphilosophique}
	\stepcounter{sujetreflitteraire\themeprogrammeqdeux}
	\stepcounter{nosujetreflitteraire}
	}
	  \IfStrEqCase{\lieu}{%
    {Métropole}{\def\centexam{met}}%
    {Amérique du Nord}{\def\centexam{amnord}}%
    {Amérique du Sud}{\def\centexam{amsud}}%
    {Centres étrangers}{\def\centexam{cea}}%
    {Liban}{\def\centexam{lib}}%
    {Polynésie}{\def\centexam{pol}}%
    {Asie}{\def\centexam{asi}}%
    {Sujet 0}{\def\centexam{suj}}%
    {Nouvelle-Calédonie}{\def\centexam{nca}}%
    {La Réunion}{\def\centexam{reu}}%
    {Antilles}{\def\centexam{ang}}
  }[\def\baseprefix{unk}]%
  \stepcounter{sujet\centexam\inter\themeprogrammequn}%
	\stepcounter{sujet\centexam\refle\themeprogrammeqdeux}%
	\stepcounter{sujet\centexam\themeprogrammequn}%
	\stepcounter{sujet\centexam\themeprogrammeqdeux}%
	\ifthenelse{\equal{\typequn}{\litt}}%
	%%%Si l'interprétation est littéraire, on augmente l'interprétation litt. liés et la réf. philosophique ; ainsi que les compteurs liés.
	{\stepcounter{sujet\centexam\inter\litte\themeprogrammequn}
	\stepcounter{nosujet\centexam\inter\litte}
	\stepcounter{sujet\centexam\refle\phil\themeprogrammeqdeux}
	\stepcounter{nosujet\centexam\refle\phil}%
	}
	{%%%Sinon, on augmente l'interprétation phil. et la réf. litt. ; ainsi que les compteurs liés.
	\stepcounter{sujet\centexam\inter\phil\themeprogrammequn}
	\stepcounter{nosujet\centexam\inter\phil}
	\stepcounter{sujet\centexam\refle\litte\themeprogrammeqdeux}
	\stepcounter{nosujet\centexam\refle\litte}
	}
}

%#1 = type : \inter/\refle #2 discipline : = \litte/\phil #3 = theme
\NewDocumentCommand{\statnat}{m m m}{%
  \ifnum\value{sujet#1#2#3}=0 % Plus aucun risque d'accent ici !
    \textcolor{black!70}{\arabic{sujet#1#2#3}}%
  \else
    % Le lien se construit tout seul avec les variables brutes
    \hyperlink{#1#2#3}{\arabic{sujet#1#2#3}}%
  \fi
}
  
\NewDocumentCommand{\totstatthexam}{O{} m m m}{
\ifnum\value{sujet#1all#2#3#4}=0
    \textcolor{gray!50}{\arabic{sujet#1all#2#3#4}}%
  \else
    \arabic{sujet#1all#2#3#4}%
  \fi}
%\newcommand{\txttotstat}[1]{sujetall#1}

\NewDocumentCommand{\statthexam}{O{} m m m}{
\ifnum\value{sujet#1#2#3#4}=0
\textcolor{gray!50}{\arabic{sujet#1#2#3#4}}%
\else
\arabic{sujet#1#2#3#4}%
\fi}

%%#1 : acronyme #2 : Nom du lieu
\NewDocumentCommand{\statexam}{O{} m}{
\gdef\centreexamen{#2}
\pagestyle{stat-centreexam}
\begin{table}[h]
	\global\setlength{\aboverulesep}{0.4pt}
%	\global\setlength{\belowrulesep}{0.6pt}
\begin{tabularx}{\linewidth}{M{2cm}*{8}{X}*{8}{X}}
\hline
\multicolumn{17}{c}{\makecell[cc]{{\large \textbf{#2\stepcounter{subsection}\phantomsection\addcontentsline{toc}{subsection}{#2}}}}}\\
\hline
{\tiny ~}&{\tiny ~}&{\tiny ~}&{\tiny ~}&{\tiny ~}&{\tiny ~}&{\tiny ~}&{\tiny ~}&{\tiny ~}&{\tiny ~}&{\tiny ~}&{\tiny ~}&{\tiny ~}&{\tiny ~}&{\tiny ~}&{\tiny ~}&{\tiny ~}\\\hline
\rowcolor{lightgray!60}
~ & \multicolumn{8}{>{\hsize=\dimexpr8\hsize+14\tabcolsep\relax\centering\arraybackslash}X}{\makecell[cc]{\textbf{La Recherche de Soi}}} 
  & \multicolumn{8}{>{\hsize=\dimexpr8\hsize+14\tabcolsep\relax\centering\arraybackslash}X}{\makecell[cc]{\textbf{L'Humanité en question}}} \\ \hline
\vspace{6pt}~&\multicolumn{3}{c}{\makecell[cc]{\textit{Exclusivement}}}&\multicolumn{5}{c}{\makecell[cc]{\textit{Plusieurs chap.}}}&\multicolumn{3}{c}{\makecell[cc]{\textit{Exclusivement}}}&\multicolumn{5}{c}{\makecell[cc]{\textit{Plusieurs chap.}}}\\\cmidrule(l){2-4}\cmidrule(l){5-9}\cmidrule(l){10-12}\cmidrule(l){13-17}
Questions & \tvert{ETA} & \tvert{EXS} & \tvert{MM} & \tvert{ETAMM} & \tvert{EXSTA} & \tvert{EXSMM} & \tvert{OdS} & \tvert{Tot.} & \tvert{CC} & \tvert{HV} & \tvert{HL} & \tvert{HVCC} & \tvert{HLCC} & \tvert{HVL} & \tvert{OdH} & \tvert{Tot.}\\\cmidrule{1-4}\cmidrule(l){5-9}\cmidrule(l){10-12}\cmidrule(l){13-17}
Int. litt. & \statthexam[#1]{\inter}{\litte}{ETA} & \statthexam[#1]{\inter}{\litte}{EXS} & \statthexam[#1]{\inter}{\litte}{MM} & \statthexam[#1]{\inter}{\litte}{ETAMM} & \statthexam[#1]{\inter}{\litte}{EXSTA} & \statthexam[#1]{\inter}{\litte}{EXSMM} & \statthexam[#1]{\inter}{\litte}{OdS} & \arabic{sujet#1RS\inter\litte} & \statthexam[#1]{\inter}{\litte}{CC} & \statthexam[#1]{\inter}{\litte}{HV} & \statthexam[#1]{\inter}{\litte}{HL} & \statthexam[#1]{\inter}{\litte}{HVCC} & \statthexam[#1]{\inter}{\litte}{HLCC} & \statthexam[#1]{\inter}{\litte}{HVL} & \statthexam[#1]{\inter}{\litte}{OdH} & \arabic{sujet#1HQ\inter\litte} \\\cmidrule(l){2-4}\cmidrule(l){5-9}\cmidrule(l){10-12}\cmidrule(l){13-17}
Int. phil. & \statthexam[#1]{\inter}{\phil}{ETA} & \statthexam[#1]{\inter}{\phil}{EXS} & \statthexam[#1]{\inter}{\phil}{MM} & \statthexam[#1]{\inter}{\phil}{ETAMM} & \statthexam[#1]{\inter}{\phil}{EXSTA} & \statthexam[#1]{\inter}{\phil}{EXSMM} & \statthexam[#1]{\inter}{\phil}{OdS} & \arabic{sujet#1RS\inter\phil} & \statthexam[#1]{\inter}{\phil}{CC} & \statthexam[#1]{\inter}{\phil}{HV} & \statthexam[#1]{\inter}{\phil}{HL} & \statthexam[#1]{\inter}{\phil}{HVCC} & \statthexam[#1]{\inter}{\phil}{HLCC} & \statthexam[#1]{\inter}{\phil}{HVL} & \statthexam[#1]{\inter}{\phil}{OdH} & \arabic{sujet#1HQ\inter\phil} \\\cmidrule(l){2-4}\cmidrule(l){5-9}\cmidrule(l){10-12}\cmidrule(l){13-17}
Essai litt. & \statthexam[#1]{\refle}{\litte}{ETA} & \statthexam[#1]{\refle}{\litte}{EXS} & \statthexam[#1]{\refle}{\litte}{MM} & \statthexam[#1]{\refle}{\litte}{ETAMM} & \statthexam[#1]{\refle}{\litte}{EXSTA} & \statthexam[#1]{\refle}{\litte}{EXSMM} & \statthexam[#1]{\refle}{\litte}{OdS} & \arabic{sujet#1RS\refle\litte} & \statthexam[#1]{\refle}{\litte}{CC} & \statthexam[#1]{\refle}{\litte}{HV} & \statthexam[#1]{\refle}{\litte}{HL} & \statthexam[#1]{\refle}{\litte}{HVCC} & \statthexam[#1]{\refle}{\litte}{HLCC} & \statthexam[#1]{\refle}{\litte}{HVL} & \statthexam[#1]{\refle}{\litte}{OdH} & \arabic{sujet#1HQ\refle\litte} \\\cmidrule(l){2-4}\cmidrule(l){5-9}\cmidrule(l){10-12}\cmidrule(l){13-17}
Essai phil. & \statthexam[#1]{\refle}{\phil}{ETA} & \statthexam[#1]{\refle}{\phil}{EXS} & \statthexam[#1]{\refle}{\phil}{MM} & \statthexam[#1]{\refle}{\phil}{ETAMM} & \statthexam[#1]{\refle}{\phil}{EXSTA} & \statthexam[#1]{\refle}{\phil}{EXSMM} & \statthexam[#1]{\refle}{\phil}{OdS} & \arabic{sujet#1RS\refle\phil} & \statthexam[#1]{\refle}{\phil}{CC} & \statthexam[#1]{\refle}{\phil}{HV} & \statthexam[#1]{\refle}{\phil}{HL} & \statthexam[#1]{\refle}{\phil}{HVCC} & \statthexam[#1]{\refle}{\phil}{HLCC} & \statthexam[#1]{\refle}{\phil}{HVL} & \statthexam[#1]{\refle}{\phil}{OdH} & \arabic{sujet#1HQ\refle\phil} \\\hline
%~&~&~&~&~&~&~&~&~&~&~&~&~&~&~&~&~\\\hline
\end{tabularx}

\vspace{6pt}


\begin{tabularx}{\linewidth}{M{2cm} *{9}{X} S *{6}{X}}
    \cmidrule(r){1-10}\cmidrule{12-17}\\*[-18.75pt]
    \rowcolor{lightgray!60}
    ~ & \multicolumn{9}{c}{\textit{Transversaux sur les semestres}} 
      & \cellcolor{white}
      & \multicolumn{6}{c}{\textit{Transversaux sur les années}} \\*[-1pt]
\cmidrule(r){1-10}\cmidrule{12-17}
& \tvert{ETACC} & \tvert{ETAHV} & \tvert{ETAHL} & \tvert{EXSCC} & \tvert{EXSHV} & \tvert{EXSHL} & \tvert{MMCC} & \tvert{MMHV} & \tvert{MMHL} &~& \tvert{ETA Prem} & \tvert{EXS Prem} & \tvert{MM Prem} & \tvert{CC Prem} & \tvert{HV Prem} & \tvert{HL Prem} \\\cmidrule(r){1-10}\cmidrule{12-17}
Int. litt. & \statthexam[#1]{\inter}{\litte}{ETACC} & \statthexam[#1]{\inter}{\litte}{ETAHV} & \statthexam[#1]{\inter}{\litte}{ETAHL} & \statthexam[#1]{\inter}{\litte}{EXSCC} & \statthexam[#1]{\inter}{\litte}{EXSHV} & \statthexam[#1]{\inter}{\litte}{EXSHL} & \statthexam[#1]{\inter}{\litte}{MMCC} & \statthexam[#1]{\inter}{\litte}{MMHV} & \statthexam[#1]{\inter}{\litte}{MMHL} &~&  \statthexam[#1]{\inter}{\litte}{premtermETA} & \statthexam[#1]{\inter}{\litte}{premtermEXS} & \statthexam[#1]{\inter}{\litte}{premtermMM} & \statthexam[#1]{\inter}{\litte}{premtermCC} & \statthexam[#1]{\inter}{\litte}{premtermHV} & \statthexam[#1]{\inter}{\litte}{premtermHL} \\\cmidrule(lr){2-10}\cmidrule{12-17}
Int. phil. & \statthexam[#1]{\inter}{\phil}{ETACC} & \statthexam[#1]{\inter}{\phil}{ETAHV} & \statthexam[#1]{\inter}{\phil}{ETAHL} & \statthexam[#1]{\inter}{\phil}{EXSCC} & \statthexam[#1]{\inter}{\phil}{EXSHV} & \statthexam[#1]{\inter}{\phil}{EXSHL} & \statthexam[#1]{\inter}{\phil}{MMCC} & \statthexam[#1]{\inter}{\phil}{MMHV} & \statthexam[#1]{\inter}{\phil}{MMHL} &~&  \statthexam[#1]{\inter}{\phil}{premtermETA} & \statthexam[#1]{\inter}{\phil}{premtermEXS} & \statthexam[#1]{\inter}{\phil}{premtermMM} & \statthexam[#1]{\inter}{\phil}{premtermCC} & \statthexam[#1]{\inter}{\phil}{premtermHV} & \statthexam[#1]{\inter}{\phil}{premtermHL} \\\cmidrule(lr){2-10}\cmidrule{12-17}
Essai litt. & \statthexam[#1]{\refle}{\litte}{ETACC} & \statthexam[#1]{\refle}{\litte}{ETAHV} & \statthexam[#1]{\refle}{\litte}{ETAHL} & \statthexam[#1]{\refle}{\litte}{EXSCC} & \statthexam[#1]{\refle}{\litte}{EXSHV} & \statthexam[#1]{\refle}{\litte}{EXSHL} & \statthexam[#1]{\refle}{\litte}{MMCC} & \statthexam[#1]{\refle}{\litte}{MMHV} & \statthexam[#1]{\refle}{\litte}{MMHL} &~&  \statthexam[#1]{\refle}{\litte}{premtermETA} & \statthexam[#1]{\refle}{\litte}{premtermEXS} & \statthexam[#1]{\refle}{\litte}{premtermMM} & \statthexam[#1]{\refle}{\litte}{premtermCC} & \statthexam[#1]{\refle}{\litte}{premtermHV} & \statthexam[#1]{\refle}{\litte}{premtermHL} \\\cmidrule(lr){2-10}\cmidrule{12-17}
Essai phil. & \statthexam[#1]{\refle}{\phil}{ETACC} & \statthexam[#1]{\refle}{\phil}{ETAHV} & \statthexam[#1]{\refle}{\phil}{ETAHL} & \statthexam[#1]{\refle}{\phil}{EXSCC} & \statthexam[#1]{\refle}{\phil}{EXSHV} & \statthexam[#1]{\refle}{\phil}{EXSHL} & \statthexam[#1]{\refle}{\phil}{MMCC} & \statthexam[#1]{\refle}{\phil}{MMHV} & \statthexam[#1]{\refle}{\phil}{MMHL} &~&  \statthexam[#1]{\refle}{\phil}{premtermETA} & \statthexam[#1]{\refle}{\phil}{premtermEXS} & \statthexam[#1]{\refle}{\phil}{premtermMM} & \statthexam[#1]{\refle}{\phil}{premtermCC} & \statthexam[#1]{\refle}{\phil}{premtermHV} & \statthexam[#1]{\refle}{\phil}{premtermHL} \\\cmidrule(lr){1-10}\cmidrule{12-17}
\end{tabularx}

\vspace*{6pt}

\begin{multicols}{2}
\begin{tabularx}{0.9\linewidth}{M{2cm}XXXXXX}
\hline
\rowcolor{lightgray!60}~&\multicolumn{6}{c}{\textit{Tous les axes confondus}}\\\hline
\vspace*{6pt}
~&~ETA~&~EXS~&~MM~&~CC~&~HV~&~HL~\\\cmidrule(l){2-7}
Int. litt.& \totstatthexam[#1]{\inter}{\litte}{ETA} & \totstatthexam[#1]{\inter}{\litte}{EXS} & \totstatthexam[#1]{\inter}{\litte}{MM} & \totstatthexam[#1]{\inter}{\litte}{CC} & \totstatthexam[#1]{\inter}{\litte}{HV} & \totstatthexam[#1]{\inter}{\litte}{HL}\\\cmidrule(l){2-7}
Int. phil.& \totstatthexam[#1]{\inter}{\phil}{ETA} & \totstatthexam[#1]{\inter}{\phil}{EXS} & \totstatthexam[#1]{\inter}{\phil}{MM} & \totstatthexam[#1]{\inter}{\phil}{CC} & \totstatthexam[#1]{\inter}{\phil}{HV} & \totstatthexam[#1]{\inter}{\phil}{HL}\\\cmidrule(l){2-7}
Essai. litt.& \totstatthexam[#1]{\inter}{\litte}{ETA} & \totstatthexam[#1]{\refle}{\litte}{EXS} & \totstatthexam[#1]{\refle}{\litte}{MM} & \totstatthexam[#1]{\refle}{\litte}{CC} & \totstatthexam[#1]{\refle}{\litte}{HV} & \totstatthexam[#1]{\refle}{\litte}{HL} \\\cmidrule(l){2-7}
Essai. phil.& \totstatthexam[#1]{\inter}{\phil}{ETA} & \totstatthexam[#1]{\refle}{\phil}{EXS} & \totstatthexam[#1]{\refle}{\phil}{MM} & \totstatthexam[#1]{\refle}{\phil}{CC} & \totstatthexam[#1]{\refle}{\phil}{HV} & \totstatthexam[#1]{\refle}{\phil}{HL}\\\cmidrule(l){2-7}
Totaux & \arabic{sujet#1allETA} & \arabic{sujet#1allEXS} & \arabic{sujet#1allMM} & \arabic{sujet#1allCC} & \arabic{sujet#1allHV} & \arabic{sujet#1allHL}\\\hline
\end{tabularx}

\begin{center}
Au total, il y a eu \number\numexpr\value{sujet#1all}/2\relax{} 
%\arabic{sujet#1all}{}
sujets proposés en #2.
\end{center}

\columnbreak
\setstretch{0.8}{{\scriptsize \underline{Liste des abréviations :} 
\underline{ETA~:} \ETA{}, \underline{EXS~:} \EXS{}, \underline{MM~:} \MM{}, \underline{CC~:} \CC, \underline{HV~:} \HV{}, \underline{HL~:} \HL{}, \underline{ETAMM} \ETAMM, \underline{EXSTA~:} \EXS{} et \ETA{}, \underline{HVL~:} \HVL, \underline{EXSMM~:} \EXSMM, \underline{ETACC~:} \ETACC, \underline{ETAHV~:} \ETAHV{}, \underline{ETAHL~:} Éducation, transmission et émancipation et \HL, \underline{EXSCC~:} \EXSCC, \underline{EXSHV~:} \EXSHV, \underline{EXSHL~:} \EXSHL, \underline{MMCC~:} \MMCC, \underline{MMHV~:} \MMHV, \underline{MMHL~:} \MMHL. \underline{« Prem » :} sujet croisé entre un chapitre de Terminale et un chapitre de Première. Section \underline{« Tous les axes confondus »~:} nombre de sujets au total qui portent sur un chapitre donné, en comptant également les sujets « mixtes ».}}
\end{multicols}
\end{table}
}%Fin de la commande statexam

%5.2 Index de l'auteur pour la table des matières
\newcommand{\indexerettoc}{%
  % Étape 1 : Format de l'auteur pour la table des matières
  \def\tocauteurformat{}%
  \ifthenelse{\equal{\prenomauteur}{}}%
    {% Sans prénom
      \def\tocauteurformat{\textsc{\nomauteur}}%
    }%
    {% Avec prénom
    \ifthenelse{\equal{\particule}{}}
    	{\def\tocauteurformat{\initialeprenom{}.~\textsc{\nomauteur}}}
    	{\IfEndWith{\particule}{'}
    	{\def\tocauteurformat{\initialeprenom{}.~\textsc{\particule\nomauteur}}}
    	{\def\tocauteurformat{\initialeprenom{}.~\textsc{\particule{}~\nomauteur}}}
    	}
    	}
  
  % Étape 2 : Format de l'auteur pour l'index
  \def\indexauteurformat{}%
  \ifthenelse{\equal{\prenomauteur}{}}%
    {% Sans prénom
    \def\indexauteurformat{\nomauteur}%
    }%
    {% Avec prénom
    \ifthenelse{\equal{\particule}{}}
    {\def\indexauteurformat{{\initialeprenom{}}.~\nomauteur}}
    {\IfEndWith{\particule}{'}
    {\def\indexauteurformat{{\initialeprenom{}}.~\particule\nomauteur}}
    {\def\indexauteurformat{{\initialeprenom{}}.~\particule{}~\nomauteur}}
    }
    }%
  
  % Étape 3 : Ajout à la table des matières
  \phantomsection
  \ifthenelse{\equal{\type}{}}%
    {% Épreuve obligatoire
      \addcontentsline{toc}{section}{\tocauteurformat, \textit{\titre} (\lieu~\annee, Jour~\jour)}%
    }%
    {% Épreuve avec type spécifié
      \addcontentsline{toc}{section}{\tocauteurformat, \textit{\titre} (\lieu~\annee~(\type), Jour~\jour)}%
    }%

  % Étape 4 : Index
\ifthenelse{\equal{\type}{}}%
      {\index[n]{\annee~\lieu~(Jour~\jour) -- \nomauteur{}@\textbf{\annee,~\lieu} (Jour~\jour) -- \\*\hspace*{1em}\indexauteurformat, \textit{\titre}}}% Épreuve obligatoire
      {\index[n]{\annee~\lieu~(\type, Jour~\jour) -- \nomauteur{}@\textbf{\annee,~\lieu} (\type, Jour~\jour) -- \\*\hspace*{1em}\indexauteurformat, \textit{\titre}}}% Épreuve avec type spécifié
      }
%%%%%%Génération automatique des liens vers les corrigés%%%%%%%%%%%

\newcommand{\resetliencorrige}{%
  \expandafter\def\csname cmdSUJ@sujet@liencorrige\endcsname{}%
}

\newcommand{\createlink}{%
  % Forcer l'expansion de \liencorrige pour le test
  \edef\templien{\liencorrige}%
  \ifthenelse{\equal{\templien}{}}{%
  % Préfixe selon le lieu
  \def\baseprefix{}%
  \IfStrEqCase{\lieu}{%
    {Métropole}{\def\baseprefix{met}}%
    {Amérique du Nord}{\def\baseprefix{amnord}}%
    {Amérique du Sud}{\def\baseprefix{amsud}}%
    {Antilles}{\def\baseprefix{ang}}%
    {Centres étrangers}{\def\baseprefix{cea}}%
    {Liban}{\def\baseprefix{lib}}%
    {Polynésie}{\def\baseprefix{pol}}%
    {Asie}{\def\baseprefix{asi}}%
    {Sujet 0}{\def\baseprefix{suj}}%
    {Nouvelle-Calédonie}{\def\baseprefix{nca}}%
    {La Réunion}{\def\baseprefix{reu}}%
  }[\def\baseprefix{unk}]%
  
  % Suffixe selon la nature de l'épreuve
  \def\suffixe{}%
  \IfStrEq{\type}{Remplacement}{\def\suffixe{rem}}{}%
  \IfStrEq{\type}{Rattrapage}{\def\suffixe{rat}}{}%
%  \ifStrEq{\type}{Candidats libres}{\def\suffixe{cal}}{}%
  
  % Utiliser le compteur nosujet qui existe déjà
%  \def\numeroSujet{\arabic{nosujet}}% 
  % Construction du lien
 \expandafter\edef\csname cmdSUJ@sujet@liencorrige\endcsname{\baseprefix\suffixe_\annee_\jour}
  }{}}
  
%%%%%%%Gestion de la base de données csv%%%%%%%%%%%%
\newcommand{\creationcsv}{%
\ifgencsv
\newwrite\sujetdatabase
\immediate\openout\sujetdatabase=sujetsdata.csv
\immediate\write\sujetdatabase{%
  numero;lieu;annee;jour;nomauteur;titre;typequn;axeun;typeqdeux;axedeux;%
  liencorrige;questun;questdeux%
}%
\fi
}

\newcommand{\writedonneecsv}{%
  \ifgencsv
	\createlink
  \def\axeun{\themeprogrammequn}%
  \def\axedeux{\themeprogrammeqdeux}%
  \immediate\write\sujetdatabase{%
    \arabic{nosujet};\lieu;\annee;\jour;%
    \nomauteur;\titre;%
    \typequn;\axeun;\typeqdeux;\axedeux;%
    \liencorrige;%
     \detokenize\expandafter{\cmdSUJ@sujet@qun}; \detokenize\expandafter{\cmdSUJ@sujet@qdeux}%
  }%
  \fi
}

\ifgencsv
\creationcsv
\fi
\newcommand{\sujetn}[2]{
\stepcounter{nosujet}
%%%%%Pour la mise en page%%%%
\renewcommand{\headrulewidth}{0.4pt}
\fancyhead{}
\fancyfoot{} %clear all footer fields
\fancyfoot[RO]{\textbf{\thepage ~sur~\pageref*{TotPages}}\hspace*{1.5cm}}
\fancyfoot[LE]{\hspace*{1.5cm}\textbf{\thepage ~sur~\pageref*{TotPages}}}
  \fancyfoot[LO,RE]{\scriptsize{Humanités, Littérature, Philosophie -- Banque des sujets indexés -- Terminale}}
  %%%Réinitialse liencorrige, sinon la création automatique de lien ne fonctionne pas.
  \resetliencorrige%
    \setkeys[SUJ]{sujet}{#1}%
    \createlink%%Création du lien vers le sujet et vers le corrigé s'il n'existe pas.
\ifgencsv
\writedonneecsv
\fi
     % En-tête avec le formatage du titre
	\fancyhead[C]{\formatsessionentete{}\formatauteur{}, \textit{\titre}}
	\indexerettoc{}
\renewcommand{\arraystretch}{1.5}
\vspace{-0.5cm}
\begin{center}
\begin{huge}
\textit{Sujet d'Humanités, Littérature, Philosophie}
\end{huge}\\
{\Large{Épreuve de Terminale}}\par
\smallskip
\hspace*{0.7cm}\noindent\fbox{%%Mise en page normal dans un petit encart au début de la page
\makebox[\linewidth-3.25cm][c]{
\begin{minipage}{\linewidth-3.25cm}
		\begin{small}
		\noindent\small \textbf{Axe du programme~:}~\semestre \hfill{
		\noindent\textbf{Genre du texte~:}~\genre}\par 
		\noindent\textbf{Auteur~:}
		\ifthenelse{\equal{\prenomauteur}{}}%%%Si l'auteur n'a pas de prénom
		{\textsc{\nomauteur}}%
		%%%Si l'auteur a un prénom
		{\ifthenelse{\equal{\particule}{}}
		{\initialeprenom{}.~\textsc{\nomauteur}}
		{\IfEndWith{\particule}{'}
		{\initialeprenom{}.~\textsc{\particule\nomauteur}}
		{\initialeprenom{}.~\textsc{\particule}~\textsc{\nomauteur}}
		}
		}%%%Fin du cas où l'auteur a un prénom
		\hfill \textbf{Ouvrage~:} \textit{\titre}\par
		\textbf{Centre d'examen :} \lieu \hfill \textbf{Année :} \annee\par
		\textbf{Session~:~}%
		\ifthenelse{\equal{\type}{}}{Obligatoire, Jour~\jour}{\type, Jour~\jour}%, Sujet \no{}\arabic{nosujet}%Pour vérifier que le nombre de sujets est bien conforme
		\ifthenelse{\equal{\corrige}{}}{\hfill{}\textbf{Éléments d'évaluation disponibles~:} Non\hypertarget{\liencorrige}{}}{\hfill{}\textbf{Éléments d'évaluation disponibles~:} {\hyperlink{corr\liencorrige}{Oui}\hypertarget{\liencorrige}{}}}\end{small}
\end{minipage}
	\indexerauteur{}%
	}
	}%%%Fin de la fbox
	\newline \smallskip \newline \noindent\textbf{Le candidat traite les deux parties sur des copies séparées.}
	\end{center}\vspace{-0.3cm}
	\begin{envtexte}
	#2\par
	\vspace{-0.4cm}\begin{flushright}\nopagebreak
	\ifthenelse{\equal{\prenomauteur}{}}
	{\textsc{\nomauteur}}%L'auteur n'a pas de prénom
	{\prenomauteur~\ifthenelse{\equal{\particule}{}}
	{\textsc{\nomauteur}}%L'auteur a un nom sans particule
	{\IfEndWith{\particule}{'}
	{\textsc{\particule\nomauteur}}%La particule s'achève par '
	{\textsc{\particule{}~\nomauteur}}%Particule ordinaire
	}%Fin du cas à particule
	}, \textit{\titre} (\datepub)%
	\ifthenelse{\equal{\ajoutref}{}}{.}{, \ajoutref{}.}
	\end{flushright}
	\end{envtexte}
 \indexerquestions{}%
 \compteurquestions{}%
	\clearpage
    }
%%%%%%% Gestion des graphiques par centre d'examen%%%%

%% Transformation de la valeur en donnée numérique gérée avec pgf-pie
\newcommand{\getcompteur}[1]{\number\value{#1}}
%\newcommand{\getcompteurs}[2]{\number\value{#1#2}}

\NewDocumentCommand{\camembertthemesgen }{O{} m}{% #1 = centre (ex. asi, met, amnord); #2 = Titre
  \edef\totalsujetscentre{\number\value{sujet#1total}}%
  \edef\pctETA{\fpeval{round(100 * \value{sujet#1allETA} / \totalsujetscentre, 0)}}%
  \edef\pctEXS{\fpeval{round(100 * \value{sujet#1allEXS} / \totalsujetscentre, 0)}}%
  \edef\pctMM{\fpeval{round(100 * \value{sujet#1allMM} / \totalsujetscentre, 0)}}%
  \edef\pctCC{\fpeval{round(100 * \value{sujet#1allCC} / \totalsujetscentre, 0)}}%
  \edef\pctHV{\fpeval{round(100 * \value{sujet#1allHV} / \totalsujetscentre, 0)}}%
  \edef\pctHL{\fpeval{round(100 * \value{sujet#1allHL} / \totalsujetscentre, 0)}}%
  \edef\pctaut{\fpeval{round(100 * \value{sujet#1allaut} / \totalsujetscentre, 0)}}%
  \begin{center}
  \textbf{#2}\\
  \begin{tikzpicture}
    \pie[radius=2, text=legend,
         sum=\totalsujetscentre,
         hide number=true,
         color={purple!50,red!60,orange!50,yellow!50,green!50,blue!50,gray!50}]{%
      \getcompteur{sujet#1allETA}/ETA (\pctETA\%),
      \getcompteur{sujet#1allEXS}/EXS (\pctEXS\%),
      \getcompteur{sujet#1allMM}/MM (\pctMM\%),
      \getcompteur{sujet#1allCC}/CC (\pctCC\%),
      \getcompteur{sujet#1allHV}/HV (\pctHV\%),
      \getcompteur{sujet#1allHL}/HL (\pctHL\%),
      \getcompteur{sujet#1allaut}/Aut. (\pctaut\%)%
    }%
  \end{tikzpicture}%
  \end{center}
}

\NewDocumentCommand{\camembertthemesexcl}{O{} m}{%
  \edef\totalsujetscentre{\number\value{sujet#1all}}%
  \edef\pctexclRS{\fpeval{round(100 * \value{sujet#1exclRS} / \totalsujetscentre, 0)}}%
  \edef\pctmixtRS{\fpeval{round(100 * \value{sujet#1mixtRS} / \totalsujetscentre, 0)}}%
  \edef\pctexclHQ{\fpeval{round(100 * \value{sujet#1exclHQ} / \totalsujetscentre, 0)}}%
  \edef\pctmixtHQ{\fpeval{round(100 * \value{sujet#1mixtHQ} / \totalsujetscentre, 0)}}%
  \edef\pctmixtHQRS{\fpeval{round(100 * \value{sujet#1mixtHQRS} / \totalsujetscentre, 0)}}%
  \begin{center}
    \textbf{#2}\\
    \begin{tikzpicture}
      \pie[radius=2, text=legend,
           sum=\totalsujetscentre,
           color={red!55,red!45,blue!60,blue!50,yellow!60}]{%
            \getcompteur{sujet#1exclRS}/Excl. RS (\pctexclRS\%),
            \getcompteur{sujet#1mixtRS}/Mixt. RS (\pctmixtRS\%),
            \getcompteur{sujet#1exclHQ}/Excl. HQ (\pctexclHQ\%),
            \getcompteur{sujet#1mixtHQ}/Mixt. HQ (\pctmixtHQ\%),
            \getcompteur{sujet#1mixtHQRS}/Mixt. HQRS (\pctmixtHQRS\%)%
      }%
    \end{tikzpicture}%
  \end{center}
}

\NewDocumentCommand{\camembertthemes}{O{} m m}{% #1 = centre (ex. asi, met, amnord); #2 = discipline; #3 = Titre
  \edef\totalsujetscentre{\number\value{sujet#1#2total}}%
  \edef\pctETA{\fpeval{round(100 * \value{sujet#1#2ETA} / \totalsujetscentre, 0)}}%
  \edef\pctEXS{\fpeval{round(100 * \value{sujet#1#2EXS} / \totalsujetscentre, 0)}}%
  \edef\pctMM{\fpeval{round(100 * \value{sujet#1#2MM} / \totalsujetscentre, 0)}}%
  \edef\pctCC{\fpeval{round(100 * \value{sujet#1#2CC} / \totalsujetscentre, 0)}}%
  \edef\pctHV{\fpeval{round(100 * \value{sujet#1#2HV} / \totalsujetscentre, 0)}}%
  \edef\pctHL{\fpeval{round(100 * \value{sujet#1#2HL} / \totalsujetscentre, 0)}}%
  \edef\pctaut{\fpeval{round(100 * \value{sujet#1#2aut} / \totalsujetscentre, 0)}}%
  \begin{center}
  \textbf{#3}\\
  \begin{tikzpicture}
    \pie[radius=2, text=legend,
         sum=\totalsujetscentre,
         hide number=true,
         color={purple!50,red!60,orange!50,yellow!50,green!50,blue!50,gray!50}]{%
      \getcompteur{sujet#1#2ETA}/ETA (\pctETA\%),
      \getcompteur{sujet#1#2EXS}/EXS (\pctEXS\%),
      \getcompteur{sujet#1#2MM}/MM (\pctMM\%),
      \getcompteur{sujet#1#2CC}/CC (\pctCC\%),
      \getcompteur{sujet#1#2HV}/HV (\pctHV\%),
      \getcompteur{sujet#1#2HL}/HL (\pctHL\%),
      \getcompteur{sujet#1#2aut}/Aut. (\pctaut\%)%
    }%
  \end{tikzpicture}%
  \end{center}
}

% #1 : centre (ex. met)
% #2 : type d'épreuve (int ou ref)
% #3 : discipline (litteraire ou philosophique)
\NewDocumentCommand{\barreexam}{O{} m m}{%
    % 1. Extraction et "gel" des valeurs des compteurs
    \edef\vETA{\number\value{sujet#1#2#3ETA}}%
    \edef\vEXS{\number\value{sujet#1#2#3EXS}}%
    \edef\vMM{\number\value{sujet#1#2#3MM}}%
    \edef\vCC{\number\value{sujet#1#2#3CC}}%
    \edef\vHV{\number\value{sujet#1#2#3HV}}%
    \edef\vHL{\number\value{sujet#1#2#3HL}}%
    \edef\vEXSTA{\number\value{sujet#1#2#3EXSTA}}%
    \edef\vETAMM{\number\value{sujet#1#2#3ETAMM}}%
    \edef\vEXSMM{\number\value{sujet#1#2#3EXSMM}}%
    \edef\vHVCC{\number\value{sujet#1#2#3HVCC}}%
    \edef\vHLCC{\number\value{sujet#1#2#3HLCC}}%
    \edef\vHVL{\number\value{sujet#1#2#3HVL}}%
    % Cas particulier pour le transversal
    \edef\vTrans{\number\value{sujettrans#1#2#3}}%
    
    % 2. Calcul du plafond dynamique
    \edef\maxval{\fpeval{max(\vETA, \vEXS, \vMM, \vCC, \vHV, \vHL, \vEXSTA, \vETAMM, \vEXSMM, \vHVCC, \vHLCC, \vHVL, \vTrans) + 1}}%

    % 3. Génération du graphique
    \begin{tikzpicture}
    \begin{axis}[
        ybar, 
        bar width=6pt,
        ymin=0, ymax=\maxval, % À ajuster selon le volume maximum de ta banque
        symbolic x coords={ETA,EXS,MM,CC,HV,HL,EXSTA,ETAMM,EXSMM,HVCC,HLCC,HVL,Trans},
        xtick=data,
        x tick label style={rotate=45, anchor=east, font=\tiny},
        width=8.5cm, height=3.5cm,
        axis lines*=left,
        nodes near coords, 
        every node near coord/.append style={font=\tiny},
        % Esthétique fidèle à ta capture d'écran
        fill=black!10, draw=black
    ]
    % PGFPlots reçoit ici des nombres purs, ce qui évite les plantages
    \addplot coordinates {
        (ETA,\vETA) (EXS,\vEXS) (MM,\vMM)
        (CC,\vCC) (HV,\vHV) (HL,\vHL)
        (EXSTA,\vEXSTA) (ETAMM,\vETAMM) (EXSMM,\vEXSMM)
        (HVCC,\vHVCC) (HLCC,\vHLCC) (HVL,\vHVL)
        (Trans,\vTrans)
    };
    \end{axis}
    \end{tikzpicture}%
}

\NewDocumentCommand{\graphexamen}{O{} m}{
\clearpage\newpage
\noindent\vspace*{-6pt}
% ==========================================
% COLONNE DE GAUCHE : GÉNÉRAL
% ==========================================
\begin{minipage}[t]{0.5\textwidth}
\vspace{0pt}
    % Un cadre englobant pour la gauche, inspiré de ta capture
    \begin{tcolorbox}[
    	height=0.9835\textheight,
        colback=gray!5, 
        colframe=gray!50, 
        sharp corners, 
        boxrule=0.5pt,
        % On peut forcer la hauteur ici si on veut qu'elle s'aligne 
        % parfaitement avec le bas de "Philosophie" plus tard
    ]
        \subsection*{Graphiques pour #2}
        \camembertthemesexcl[#1]{Proportion des questions\\mixtes et exclusives\\~}
        
		\begin{center}
		{\scriptsize{\underline{Légende~:} \emph{Excl. RS} : Un seul chapitre de la Recherche de soi ; \emph{Mixt. RS} : Plusieurs chapitres de la Recherche de soi ; \emph{Excl. HQ} : Un seul chapitre de l'Humanité en question ; \emph{Mixt. HQ} : Plusieurs chapitres de l'Humanité en question ; \emph{Mixt. HQRS} : Chapitres transversaux sur les semestres et les années.}}
		\end{center}
%        \vspace{2cm} % Espacement vertical généreux entre les deux camemberts
        
        \camembertthemesgen[#1]{Proportion des questions selon\\les thèmes du programme}
        
        {\scriptsize{Dès qu'une question fait appel à un thème, y compris pour les sujets croisés, elle est comptée dans le thème. Par exemple si une question porte sur EXS et MM, elle augmentera la proportion pour EXS et pour MM.}}
    \end{tcolorbox}
\end{minipage}%
\hfill% Le ressort horizontal qui pousse la colonne de droite
% ==========================================
% COLONNE DE DROITE : Littérature & Philosophie
% ==========================================
\begin{minipage}[t]{0.85\textwidth}
    \vspace{0pt}
    % --- 1. BLOC LITTÉRATURE ---
    \begin{tcolorbox}[
    	height=0.49\textheight,
        colback=red!15, 
        colframe=black, 
        sharp corners, 
        boxrule=0.5pt
    ]
        % Séparation interne du bloc Littérature
        \begin{minipage}{0.45\textwidth}
            \begin{center} % Pour bien centrer le camembert dans sa moitié
                \camembertthemes[#1]{litteraire}{Proportion des questions selon\\les thèmes du programme}
            \end{center}
        \end{minipage}%
        \hfill%
        \begin{minipage}{0.5\textwidth}
            \textbf{Interprétation}\\
            \barreexam[#1]{int}{litteraire}
            
            \vspace{0.1cm} % Ajustement entre les deux graphiques en barres
            
            \textbf{Essai}\\
            \barreexam[#1]{ref}{litteraire}
        \end{minipage}
        
        \vspace{-0.8cm} % Espace avant d'écrire le label
        
        % Le label placé manuellement en bas à gauche
        {\large \textbf{Littérature}} 
    \end{tcolorbox}
%    
    % --- 2. BLOC PHILOSOPHIE ---
    \begin{tcolorbox}[
    	height=0.48\textheight,
        colback=gray!20, 
        colframe=black, 
        sharp corners, 
        boxrule=0.5pt
    ]
                % Séparation interne du bloc Littérature
        \begin{minipage}{0.45\textwidth}
            \begin{center} % Pour bien centrer le camembert dans sa moitié
                \camembertthemes[#1]{philosophique}{Proportion des questions selon\\les thèmes du programme}
            \end{center}
        \end{minipage}%
        \hfill%
        \begin{minipage}{0.5\textwidth}
            \textbf{Interprétation}\\
            \barreexam[#1]{int}{philosophique}
            
            \vspace{0.1cm} % Ajustement entre les deux graphiques en barres
            
            \textbf{Essai}\\
            \barreexam[#1]{ref}{philosophique}
        \end{minipage}
        
        \vspace{-1cm} % Espace avant d'écrire le label
        
        % Le label placé manuellement en bas à gauche
        {\large \textbf{Philosophie}}
    \end{tcolorbox}
\end{minipage}
}